% =============================================================
%  Dynaswarm — Algorithms in Scientific Notation
%  Part 1/3: Preamble, Introduction, Environment, A*
% =============================================================
\documentclass[11pt,a4paper]{article}
\usepackage[T1]{fontenc}
\usepackage[utf8]{inputenc}
\usepackage{lmodern}
\usepackage{amsmath,amssymb,amsthm}
\usepackage{mathtools}
\usepackage{program}          % pseudocode — available in base texlive
\usepackage{bm}
\usepackage{booktabs}
\usepackage{hyperref}
\usepackage{xcolor}
\usepackage{geometry}
\usepackage{enumitem}
\usepackage{parskip}
\geometry{margin=2.5cm}

\theoremstyle{definition}
\newtheorem{definition}{Definition}[section]
\newtheorem{proposition}{Proposition}[section]
\newtheorem{remark}{Remark}[section]

% --- notation -----------------------------------------------
\newcommand{\calG}{\mathcal{G}}
\newcommand{\calE}{\mathcal{E}}
\newcommand{\calV}{\mathcal{V}}
\newcommand{\RR}{\mathbb{R}}
\newcommand{\xvec}{\mathbf{x}}
\newcommand{\vvec}{\mathbf{v}}
\newcommand{\pvec}{\mathbf{p}}
\newcommand{\gvec}{\mathbf{g}}
\newcommand{\excess}{\varepsilon}
% =============================================================
\begin{document}

\begin{titlepage}
  \centering\vspace*{3cm}
  {\Huge\bfseries Dynaswarm}\\[0.6em]
  {\Large\itshape Adaptive Swarm Intelligence for Dynamic Path Optimisation}\\[0.4em]
  {\large A Technical Reference: Algorithms in Mathematical Notation}\\[3cm]
  {\large Vaibhav Bansal}\\[0.5em]
  {\normalsize ML Research Project --- \today}\\[4cm]
  \begin{abstract}
    \noindent
    We study whether an online adaptive controller that monitors environment volatility
    and incrementally retunes optimiser parameters can outperform a carefully tuned
    \emph{static} configuration when searching for shortest paths in dynamically disrupted
    networks.  Four algorithms are investigated: the exact reference A$^*$, the MAX--MIN
    Ant Colony Optimiser (MMAS), a priority-vector Particle Swarm Optimiser (PSO), and a
    budget-sharing hybrid.  The hybrid rests on a falsifiable prediction: the memoryless
    swarm covers the colony's staleness window immediately after a disruption.  Across 400
    independent trials the adaptive controller yields no statistically significant
    improvement; we discuss why and what the result implies for future work.
  \end{abstract}
\end{titlepage}

\tableofcontents\newpage

% =============================================================
\section{Introduction}
% =============================================================

\subsection{Motivation}

Real-world routing networks are not static. Roads close, link capacities fluctuate, and
edges are added or removed.  Two complementary strategies handle this:

\begin{enumerate}[label=(\roman*)]
  \item \textbf{Exact recomputation.}  Re-run Dijkstra or A$^*$ after every change.
        Guaranteed optimal; cost grows with graph size.
  \item \textbf{Metaheuristic search.}  Maintain a population that continuously samples
        the search space.  Sub-optimal but cheap; exploits information from prior timesteps.
\end{enumerate}

The metaheuristic approach introduces a new problem: parameters governing the
exploration--exploitation balance have an optimal value that depends on how fast the
environment changes.  The question is whether a controller that measures volatility and
adjusts these parameters in real time can outperform a fixed, well-tuned setting.

\subsection{Problem Formulation}

\begin{definition}[Dynamic Weighted Graph]
  A \emph{dynamic weighted graph} is a sequence $\{\calG^{(t)}\}_{t=0}^{T}$ where each
  $\calG^{(t)} = (\calV,\calE^{(t)},w^{(t)})$ shares a fixed vertex set
  $\calV = \{1,\ldots,n\}$, has active edge set
  $\calE^{(t)} \subseteq \calE_{\text{all}}$, and weight function
  $w^{(t)}:\calE_{\text{all}} \to \RR_{>0}$.
\end{definition}

\begin{definition}[Dynamic Shortest Path Problem]
  At each timestep $t$, find a simple path $\pi^{(t)} = (s=v_0,v_1,\ldots,v_k=q)$
  minimising total weight
  $c(\pi^{(t)}) = \sum_{i=0}^{k-1}w^{(t)}(v_i,v_{i+1})$
  subject to $(v_i,v_{i+1})\in\calE^{(t)}$ for all $i$.
\end{definition}

\begin{definition}[Excess Over Optimum]
  Let $c^*(\calG^{(t)})$ be the true optimal cost.  The excess of solution
  $\hat\pi^{(t)}$ is
  \[
    \excess^{(t)} = \frac{c(\hat\pi^{(t)})}{c^*(\calG^{(t)})} - 1 \;\ge 0.
  \]
\end{definition}

% =============================================================
\section{Environment Model}
\label{sec:environment}
% =============================================================

\subsection{Graph Representation}

The network uses a \emph{static topology superset}
$\calE_{\text{all}} = \calE_{\text{base}} \cup \calE_{\text{pool}}$
allocated once per trial.  Each undirected edge $\{u,v\}$ occupies two directed
\emph{slots}: slot $s$ for arc $(u,v)$ and slot $\texttt{rev}[s]$ for $(v,u)$.
State at time $t$:
\begin{align}
  w^{(t)}_s &\in \RR_{>0} &\text{(slot weight)}\\
  a^{(t)}_s &\in \{0,1\} &\text{(active mask)}
\end{align}
Both arrays are read-only, shared identically by every arm.

\subsection{Graph Models}

\begin{table}[h]\centering
  \caption{Graph models and their properties.}
  \begin{tabular}{llll}
    \toprule
    Model & Edge rule & Coordinates & A$^*$ heuristic \\
    \midrule
    Erd\H{o}s--R\'enyi & Each pair w.p.\ $p$ & None & Zero (Dijkstra) \\
    Geometric & Connect if $\|u-v\|_2\le r$ & 2-D unit square & Euclidean \\
    Grid & 4-connected lattice & Integer grid & Manhattan \\
    \bottomrule
  \end{tabular}
\end{table}

Edge weights for spatial models:
$w(\{u,v\}) = w_{\min} + d(u,v)(w_{\max}-w_{\min})$
where $d$ is Euclidean or Manhattan distance.

\subsection{Disruption Events}

At each timestep $t$:
\[
  K^{(t)} \sim \mathrm{Poisson}(\lambda),\quad \lambda\in\{0,0.1,0.3,0.8\}
\]
The rate is \emph{absolute} (not proportional to $|\calE|$), preserving quiet stretches.

\begin{table}[h]\centering
  \caption{Event kinds and their effect on the graph state.}
  \label{tab:events}
  \begin{tabular}{lll}
    \toprule
    Kind & Weight update & Active update \\
    \midrule
    Closure     & $w_s \leftarrow \infty$ & unchanged \\
    Traffic up  & $w_s \leftarrow \phi\,w_s$ & unchanged \\
    Traffic down& $w_s \leftarrow \max(w_s/\phi,\,w^{\text{free}}_s)$ & unchanged \\
    Reopen      & $w_s \sim \text{Uniform}[w^{\text{free}}_s,\;\phi w^{\text{free}}_s]$ & unchanged \\
    Delete      & unchanged & $a_s\leftarrow 0$ \\
    Add         & $w_s\leftarrow w^{\text{free}}_s$ & $a_s\leftarrow 1$ \\
    \bottomrule
  \end{tabular}
\end{table}

Events disconnecting $s$ from $q$ are rejected.  Cluster size: $C=25$ edges per
disruption.

\subsection{Pre-materialised Timeline}

The entire event sequence is frozen before any optimiser exists:
\[
  \mathcal{S} = (E^{(0)},\ldots,E^{(T-1)}),\quad
  \calG^{(t)} = \Phi(\calG^{(0)},E^{(0)},\ldots,E^{(t-1)}).
\]
Every arm sees byte-identical graph states verified by a BLAKE2b digest.

% =============================================================
\section{A$^*$ --- The Exact Reference}
\label{sec:astar}
% =============================================================

\subsection{Role}

A$^*$ provides $c^*(\calG^{(t)})$ at every timestep and never competes under a budget
constraint.  Its node expansion count measures the cost of exact recomputation.

\subsection{Algorithm}

$f$-score: $f(v) = g(v) + h(v)$ where $g(v)$ is the known shortest distance from $s$.

\begin{center}
\begin{minipage}{0.88\linewidth}
\hrule\medskip
\textbf{Algorithm 1:} A$^*$ shortest path $(s\to q)$

\smallskip\textit{Input:} $\calG^{(t)}$, source $s$, target $q$, heuristic $h$.
\textit{Output:} optimal path $\pi^*$, cost $c^*$.
\hrule\medskip
\begin{program}
g(v) := \infty \mbox{ for all } v;\; g(s) := 0;\; \mathit{prev}(v) := \mathbf{nil}
Q := \{(f(s),s)\} \mbox{ (min-heap)};\; \mathit{closed} := \emptyset
|while| Q \ne \emptyset |do|
  \mbox{pop } (\_,u) := \mathrm{pop\text{-}min}(Q)
  |if| u \in \mathit{closed} \mbox{ then skip}
  \mathit{closed} := \mathit{closed} \cup \{u\}
  |if| u = q \mbox{ then break}
  |for| \mbox{each active arc } (u,v) |do|
    g' := g(u) + w(u,v)
    |if| g' < g(v) \mbox{ then}
      g(v) := g';\; \mathit{prev}(v) := u
      Q := Q \cup \{(g(v)+h(v),\,v)\}
    |fi|
  |od|
|od|
|return| \mathrm{reconstruct\text{-}path}(\mathit{prev},q),\; g(q)
\end{program}
\hrule
\end{minipage}
\end{center}

\subsection{Heuristic Admissibility}

\begin{proposition}
  For geometric and grid graphs, $h(v) = w_{\mathrm{scale}}\cdot d(v,q)$ is admissible.
\end{proposition}
\begin{proof}
  Every edge satisfies $w(\{u,v\})\ge w_{\min}+d(u,v)\cdot w_{\mathrm{scale}}$, so by
  the triangle inequality $c^*(v,q)\ge w_{\mathrm{scale}}\cdot d(v,q)=h(v)$.
\end{proof}

\subsection{Behaviour Under Disruption}

A$^*$ keeps no state.  On each \texttt{rebind}$(\calG^{(t)})$ it discards its solution
and performs a full re-search, so $\excess^{(t)}_{\mathrm{A}^*}=0$ always.

% =============================================================
%  Part 2/3: MAX-MIN ACO, PSO, Hybrid
% =============================================================

\section{MAX--MIN Ant Colony Optimisation (MMAS)}
\label{sec:aco}

\subsection{Overview}

MMAS maintains a pheromone trail $\tau_s\in[\tau_{\min},\tau_{\max}]$ on each directed
slot $s$.  Ants build paths stochastically, biased toward high-pheromone low-cost arcs.
After each iteration the trail is partially evaporated and the best path is reinforced.
Trails are clamped to $[\tau_{\min},\tau_{\max}]$, preventing premature convergence.

\subsection{Visibility}

\[
  \eta_s = \frac{1}{w_s + \gamma\cdot h_s},\quad
  h_s = w_{\mathrm{scale}}\cdot d(v_s,\,q)
\]
where $\gamma$ is the goal-weight parameter and $v_s$ is the destination of slot $s$.
Without $\gamma$, ants have no notion of the target's location and settle $\approx90\%$
above the optimum on 200-node graphs.

\subsection{Transition Probability}

\begin{definition}[ACO Transition]
  At node $u$, an ant selects next hop $v$ from unvisited neighbours
  $\mathcal{N}^+(u)$:
  \[
    P(u\to v) =
    \frac{\tau_{s(u,v)}^{\alpha}\cdot\eta_{s(u,v)}^{\beta}}
         {\displaystyle\sum_{v'\in\mathcal{N}^+(u)}
          \tau_{s(u,v')}^{\alpha}\cdot\eta_{s(u,v')}^{\beta}}
  \]
  Selection by roulette-wheel sampling: draw $r\sim U[0,\Sigma]$; take the first $v$
  whose cumulative score exceeds $r$.
\end{definition}

\subsection{Trail Update (Three Steps)}

\paragraph{Step 1 — Evaporation.}
\[
  \tau_s \;\leftarrow\; (1-\rho)\,\tau_s \qquad\forall s
\]

\paragraph{Step 2 — Reinforcement.}
Let $\pi^*_{\text{iter}}$ be the best path this iteration with cost $c^*_{\text{iter}}$.
Each slot $s$ on that path and its reverse are reinforced:
\[
  \tau_s \;\mathrel{+}=\; \frac{1}{c^*_{\text{iter}}},\quad
  \tau_{\texttt{rev}[s]} \;\mathrel{+}=\; \frac{1}{c^*_{\text{iter}}}
\]

\paragraph{Step 3 — MAX--MIN Clamping.}
Let $C_{\text{best}}$ be the best cost found so far.
\begin{align}
  \tau_{\max} &= \frac{1}{\rho\cdot C_{\text{best}}} \\
  \tau_{\min} &= \frac{\tau_{\max}}{2n}
\end{align}
Clamp: $\tau_s\leftarrow\min(\tau_{\max},\max(\tau_{\min},\tau_s))$.

\begin{remark}
  Without bounds the ratio $\tau_{\max}/\tau_{\min}$ grows without bound.  At iteration 80
  on a 30-node graph it reached $9\times10^8$; every ant followed an identical deterministic
  route.  The bounds cap the ratio at $2n$, keeping the evaporation parameter meaningful.
\end{remark}

\subsection{State Under Disruption}

Trails are kept across timesteps; inactive slots are zeroed:
$\tau_s\leftarrow 0$ if $a_s^{(t)}=0$.
The stale trail on a closed road is the phenomenon under study: the controller's job is
to accelerate forgetting (raise $\rho$) when it detects disruption.

\subsection{Full Algorithm}

\begin{center}
\begin{minipage}{0.88\linewidth}
\hrule\medskip
\textbf{Algorithm 2:} MAX--MIN ACO at timestep $t$

\smallskip\textit{Input:} $\calG^{(t)}$, trails $\{\tau_s\}$,
parameters $(\alpha,\beta,\rho,n_{\text{ants}},\gamma)$, budget $B$.
\textit{Output:} updated trails; incumbent $\hat\pi$.
\hrule\medskip
\begin{program}
\mbox{Rebind: recompute }\{\eta_s\};\mbox{ zero trails on inactive slots}
b := 0 \mbox{ (evaluations spent)}
|while| b < B |do|
  F := \emptyset \mbox{ (successful paths this iteration)}
  |for| k := 1 |to| n_{\text{ants}} |do|
    |if| b \ge B |then| \mbox{break} |fi|
    \pi_k := \mathrm{ConstructPath}(\calG^{(t)},\tau,\eta,\alpha,\beta)
    b := b + 1
    |if| \pi_k \mbox{ reaches } q |then| F := F\cup\{\pi_k\} |fi|
  |od|
  \mbox{Evaporate: } \tau_s := (1-\rho)\tau_s \mbox{ for all } s
  |if| F \ne \emptyset |then|
    \pi^*_{\text{iter}} := \arg\min_{\pi\in F} c(\pi)
    \mbox{Reinforce } \pi^*_{\text{iter}} \mbox{ (both directions)}
    |if| c(\pi^*_{\text{iter}}) < c(\hat\pi) |then| \hat\pi := \pi^*_{\text{iter}} |fi|
  |fi|
  \mbox{Clamp: } \tau_s := \min(\tau_{\max},\max(\tau_{\min},\tau_s))
|od|
|return| \hat\pi,\;\{\tau_s\}
\end{program}
\hrule\medskip
\textbf{Algorithm 3:} ConstructPath (one ant)

\smallskip\textit{Input:} $\calG^{(t)}$, $\tau$, $\eta$, $\alpha$, $\beta$.
\textit{Output:} path $\pi$, or failure.
\hrule\medskip
\begin{program}
\pi := [s];\; \mathit{visited} := \{s\};\; u := s
|while| u \ne q |do|
  N := \{v \mid (u,v)\in\calE^{(t)},\; v\notin \mathit{visited}\}
  |if| N = \emptyset |then| |return| \mathrm{failure} |fi|
  \Sigma := \textstyle\sum_{v\in N}\tau_{s(u,v)}^\alpha\eta_{s(u,v)}^\beta
  \mbox{Draw } r\sim U[0,\Sigma];\; \mathit{acc} := 0
  |for| \mbox{each } v\in N |do|
    \mathit{acc} := \mathit{acc} + \tau_{s(u,v)}^\alpha\eta_{s(u,v)}^\beta
    |if| \mathit{acc} \ge r |then| u^* := v;\; \mbox{break} |fi|
  |od|
  \pi := \pi\cdot[u^*];\; \mathit{visited} := \mathit{visited}\cup\{u^*\};\; u := u^*
|od|
|return| \pi
\end{program}
\hrule
\end{minipage}
\end{center}

\subsection{Parameters}
\begin{table}[h]\centering
  \caption{MMAS tunable parameters.}
  \begin{tabular}{lllp{5.5cm}}
    \toprule Symbol & Default & Range & Interpretation \\ \midrule
    $\alpha$ & 1.0 & $[0,5]$ & Trail exponent (weight of accumulated memory) \\
    $\beta$  & 2.0 & $[0,8]$ & Visibility exponent (weight of edge cost) \\
    $\rho$   & 0.10 & $[0.01,0.90]$ & Evaporation rate (speed of forgetting) \\
    $n_{\text{ants}}$ & 20 & $[4,100]$ & Paths built per iteration \\
    $\gamma$ & 1.0 & $[0,5]$ & Goal-direction weight in visibility \\
    \bottomrule
  \end{tabular}
\end{table}

\subsection{Lambda Branching (Convergence Telemetry)}

\[
  \lambda_B = \frac{1}{|\mathcal{R}|}\sum_{u\in\mathcal{R}}
  \Bigl|\bigl\{s(u,v)\;\big|\;v\in\mathcal{N}^+(u),\;
  \tau_{s(u,v)}\ge\tau^{(u)}_{\min}+0.05(\tau^{(u)}_{\max}-\tau^{(u)}_{\min})\bigr\}\Bigr|
\]
$\mathcal{R}$: nodes recently visited.
$\lambda_B\approx 1$ means convergence; $\lambda_B\approx\bar{d}$ (mean degree) means
healthy exploration.

% =============================================================
\section{Priority-Vector Particle Swarm Optimisation}
\label{sec:pso}
% =============================================================

\subsection{Encoding}

Each particle $i$ holds a priority vector $\xvec_i\in[0,1]^n$.
Path construction: starting at $s$, choose
\[
  v^* = \arg\max_{v\in\mathcal{N}^+(u),\;v\notin\mathit{visited}}
        \xvec_i[v]\cdot\eta_{s(u,v)},\quad u\leftarrow v^*
\]
The same visibility $\eta$ as MMAS makes the two methods directly comparable.
Dead-ended particles are \emph{repaired}: priority vector is perturbed and construction
retried up to $R$ times, each attempt charged to the budget.

\begin{remark}
  A natural alternative gives each particle a bias $\mathbf{b}\in\RR^n$ and finds the
  shortest path under $w'(u,v)=w(u,v)(1+b_v)$.  This always returns a valid path, but
  the optimal solution sits at $\mathbf{b}=\mathbf{0}$ --- a distinguished point --- making
  the method equivalent to Dijkstra and providing no insight into swarm dynamics.
\end{remark}

\subsection{Velocity and Position Update}

\begin{align}
  \vvec_i^{(k+1)} &= w_I\vvec_i^{(k)}
    + c_1 r_1(\pvec_i - \xvec_i^{(k)})
    + c_2 r_2(\gvec - \xvec_i^{(k)}) \label{eq:vel}\\
  \xvec_i^{(k+1)} &= \operatorname{clip}(\xvec_i^{(k)}+\vvec_i^{(k+1)},\;0,\;1)
  \label{eq:pos}
\end{align}
where $r_1,r_2\sim U[0,1]$; $\pvec_i$ = personal best; $\gvec$ = global best.

\subsection{Handling Disruptions}

\begin{itemize}
  \item \textbf{Positions kept.}  Priorities rank nodes, not edges; a closure does not
        invalidate the ranking.
  \item \textbf{Personal costs invalidated:} $c_i\leftarrow\infty$.  Costs were priced
        against a graph that no longer exists.
  \item \textbf{Global best position kept.}  Discarding $\gvec$ was measured to cost
        $\approx90\%$ additional excess under disruption.
\end{itemize}

\subsection{Full Algorithm}

\begin{center}
\begin{minipage}{0.88\linewidth}
\hrule\medskip
\textbf{Algorithm 4:} Priority-Vector PSO at timestep $t$

\smallskip\textit{Input:} $\calG^{(t)}$, swarm $\{(\xvec_i,\vvec_i,\pvec_i)\}$,
parameters $(w_I,c_1,c_2,R,\gamma)$, budget $B$.
\textit{Output:} updated swarm; incumbent $\hat\pi$.
\hrule\medskip
\begin{program}
\mbox{Rebind: recompute }\{\eta_s\};\; c_i := \infty;\;\mbox{re-price }\hat\pi
b := 0
|while| b < B |do|
  |for| i := 1 |to| N |do|
    |if| b \ge B |then| \mbox{break} |fi|
    \mbox{Update }\vvec_i,\xvec_i \mbox{ via Eqs.\,(\ref{eq:vel})--(\ref{eq:pos})}
    \pi_i := \mathrm{Construct}(\xvec_i,\calG^{(t)},\eta);\; b := b+1
    |if| \pi_i = \mathrm{failure} |then|
      |for| r := 1 |to| R |do|
        |if| b \ge B |then| \mbox{break} |fi|
        \mbox{Perturb }\xvec_i;\; \pi_i := \mathrm{Construct}(\xvec_i,\calG^{(t)},\eta);\; b:=b+1
        |if| \pi_i \ne \mathrm{failure} |then| \mbox{break} |fi|
      |od|
    |fi|
    |if| \pi_i \ne \mathrm{failure} \mbox{ and } c(\pi_i) < c_i |then|
      c_i := c(\pi_i);\; \pvec_i := \xvec_i
    |fi|
    |if| c(\pi_i) < c(\hat\pi) |then| \hat\pi := \pi_i;\; \gvec := \xvec_i |fi|
  |od|
|od|
|return| \hat\pi
\end{program}
\hrule
\end{minipage}
\end{center}

\subsection{Parameters}
\begin{table}[h]\centering
  \caption{PSO tunable parameters.}
  \begin{tabular}{lllp{5cm}}
    \toprule Symbol & Default & Range & Interpretation \\ \midrule
    $w_I$ & 0.7 & $[0,1]$ & Inertia (velocity momentum) \\
    $c_1$ & 1.5 & $[0,4]$ & Cognitive coefficient \\
    $c_2$ & 1.5 & $[0,4]$ & Social coefficient \\
    $N$   & 20  & $[4,100]$ & Swarm size \\
    $\gamma$ & 1.0 & $[0,5]$ & Goal-direction weight \\
    $R$   & 3   & $[0,10]$ & Max repair attempts \\
    \bottomrule
  \end{tabular}
\end{table}

% =============================================================
\section{Hybrid ACO-PSO with Budget Sharing}
\label{sec:hybrid}
% =============================================================

\subsection{Rationale}

\begin{proposition}[Complementary Failure Modes]
  \begin{enumerate}[label=(\roman*)]
    \item \textbf{ACO} fails because pheromone memory points toward routes that no longer
          exist.  Recovery time $\propto 1/\rho$.
    \item \textbf{PSO} fails because the swarm has converged; disruption invalidates
          costs but not positions, leaving particles clustered in a suboptimal region of
          $\calG^{(t)}$.
  \end{enumerate}
\end{proposition}

\begin{remark}[Testable Claim]
  If the hybrid's advantage is mechanistic, it should manifest as lower
  \emph{damage} and faster \emph{recovery} in the steps immediately after a
  disruption, while steady-state excess stays comparable to ACO alone.
\end{remark}

\subsection{Budget Splitting}

\begin{align}
  B_{\mathrm{ACO}}^{(t)} &= \max\!\left(1,\;\bigl\lfloor B^{(t)}\cdot f\bigr\rceil\right) \\
  B_{\mathrm{PSO}}^{(t)} &= \max\!\left(1,\;B^{(t)}-B_{\mathrm{ACO}}^{(t)}\right)
\end{align}
$f\in[0.05,0.95]$ is the \emph{ACO fraction}.  The parent meter is charged exactly
$B_{\mathrm{ACO}}+B_{\mathrm{PSO}}$ (actual work done), preventing budget smuggling.

\subsection{Bidirectional Exchange}

\paragraph{Colony $\to$ Swarm.}
Let $\hat\pi_A$ be the colony's incumbent.  Boost swarm priorities along its nodes:
\[
  \xvec_i[v]\leftarrow\min(\xvec_i[v]+\delta_{\text{ex}},\;1)
  \quad\forall v\in\hat\pi_A,\;\forall i,\qquad \delta_{\text{ex}}=0.25
\]

\paragraph{Swarm $\to$ Colony.}
Let $\hat\pi_P$ be the swarm's incumbent.  Deposit trail on that path:
\[
  \tau_s\;\mathrel{+}=\;\frac{1}{c(\hat\pi_P)}\quad\forall s\in\hat\pi_P
  \quad\text{(both directions via \texttt{rev})}
\]

\subsection{Full Algorithm}

\begin{center}
\begin{minipage}{0.88\linewidth}
\hrule\medskip
\textbf{Algorithm 5:} Hybrid ACO-PSO at timestep $t$

\smallskip\textit{Input:} $\calG^{(t)}$, ACO state $(\tau,\hat\pi_A)$,
PSO state $(\{\xvec_i\},\hat\pi_P)$, fraction $f$, budget $B$.
\textit{Output:} updated states; best incumbent.
\hrule\medskip
\begin{program}
B_A := \max(1,\lfloor Bf \rceil);\; B_P := \max(1, B - B_A)
\hat\pi_A := \mathrm{MMAS}(\calG^{(t)},\tau,B_A) \mbox{ (Alg.\ 2)}
\hat\pi_P := \mathrm{PSO}(\calG^{(t)},\{\xvec_i\},B_P) \mbox{ (Alg.\ 4)}
\mbox{Exchange Colony}\to\mbox{Swarm:}
|if| \hat\pi_A \ne \mathbf{nil} |then|
  |for| \mbox{each } v\in\hat\pi_A, \mbox{each } i |do|
    \xvec_i[v] := \min(\xvec_i[v]+\delta_{\text{ex}},\;1)
  |od|
|fi|
\mbox{Exchange Swarm}\to\mbox{Colony:}
|if| \hat\pi_P \ne \mathbf{nil} |then|
  |for| \mbox{each slot } s \mbox{ in } \hat\pi_P |do|
    \tau_s := \tau_s + 1/c(\hat\pi_P);\;\tau_{\texttt{rev}[s]} := \tau_{\texttt{rev}[s]} + 1/c(\hat\pi_P)
  |od|
|fi|
|return| \arg\min(c(\hat\pi_A),\,c(\hat\pi_P))
\end{program}
\hrule
\end{minipage}
\end{center}

\subsection{Parameters}
The hybrid has 12 tunables: 5 ACO (prefixed \texttt{aco\_}), 6 PSO (prefixed
\texttt{pso\_}), and:

\begin{table}[h]\centering
  \caption{Hybrid-specific tunable.}
  \begin{tabular}{lllp{6cm}}
    \toprule Symbol & Default & Range & Interpretation \\ \midrule
    $f$ & 0.5 & $[0.05,0.95]$ & ACO fraction of the shared budget \\
    \bottomrule
  \end{tabular}
\end{table}

$f$ is the natural lever for allocation control: driven by the recency signal, shift
budget toward PSO (low-$f$) right after a disruption and back toward ACO during calm.

% =============================================================
%  Part 3/3: Controller, Evaluation, Results, Conclusion
% =============================================================

% =============================================================
\section{Adaptive Controller}
\label{sec:controller}
% =============================================================

\subsection{Feedback Architecture}

The controller is a discrete-time feedback loop. At timestep $t$:
\begin{enumerate}[label=(\arabic*)]
  \item The optimiser searches under current parameters $\bm{\theta}^{(t)}$.
  \item A signal $\sigma(t)$ is computed from the step's observations.
  \item The controller proposes updated parameters $\bm{\theta}^{(t+1)}$.
  \item The optimiser applies the update before the next timestep.
\end{enumerate}

\subsection{Volatility Signals}

\subsubsection{Edge Ratio Signal}
$\sigma_{\text{ER}}(t) = \frac{1}{W} \sum_{\ell=t-W+1}^{t} K^{(\ell)}/|\calE^{(\ell)}|$.
Explore direction: $+1$.
\textit{Limitation:} Dividing absolute disruptions by $|\calE|$ makes the signal small
on large graphs, failing to distinguish scenarios at the geometric tier.

\subsubsection{Volatility Signal}
$\sigma_{\text{V}}(t) = \frac{1}{W} \sum_{\ell=t-W+1}^{t} \mathbf{1}[K^{(\ell)}>0] \in [0,1]$.
Fraction of recent steps with \emph{any} disruption. Explore direction: $+1$.
\textit{Limitation:} Saturates at 1.0 when expected disruptions $\lambda \gtrsim 1$.

\subsubsection{Recency Signal (Transient)}
Tracks time since the last disruption. Decay $\mu\in(0,1)$ (default $0.7$).
\begin{align}
  \sigma_{\text{R}}(t) &=
  \begin{cases}
    1   & \text{if } K^{(t)} > 0 \\
    \mu \cdot \sigma_{\text{R}}(t-1) & \text{otherwise}
  \end{cases}
\end{align}
Explore direction: $+1$.
This is the target for the hybrid budget split $f$, as staleness is a transient property
peaking immediately after a disruption.

\subsection{Incremental Controller}

The controller compares the signal against its \emph{own running baseline} rather than
an absolute threshold.

\begin{definition}[Pressure]
  Let $\bar\sigma^{(t)}$ be the EMA of the signal with smoothing $\alpha_b$:
  \[
    \bar\sigma^{(t)} = (1-\alpha_b)\bar\sigma^{(t-1)} + \alpha_b\sigma(t)
  \]
  Pressure: $\delta^{(t)} = \xi\cdot\sigma(t) - \bar\sigma^{(t)}$, where $\xi$ is the
  signal's explore direction.
\end{definition}

\begin{center}
\begin{minipage}{0.88\linewidth}
\hrule\medskip
\textbf{Algorithm 6:} Incremental Controller (adaptation step at time $t$)

\smallskip\textit{Input:} $\sigma(t)$, $\bm{\theta}^{(t)}$, $\bar\sigma$,
targets $\mathcal{T}$, step $\Delta$, deadband $\varepsilon_d$, period $p$.
\textit{Output:} $\bm{\theta}^{(t+1)}$, updated baseline $\bar\sigma$.
\hrule\medskip
\begin{program}
\delta := \xi \cdot \sigma(t) - \bar\sigma
\bar\sigma := \bar\sigma + \alpha_b\cdot\delta
|if| t \bmod p \ne 0 \mbox{ or } |\delta| \le \varepsilon_d |then|
  |return| \bm{\theta}^{(t)},\;\bar\sigma \mbox{ (no change)}
|fi|
\bm{\theta}^{(t+1)} := \bm{\theta}^{(t)}
|for| \mbox{each target } \theta_k\in\mathcal{T} |do|
  \mathit{step} := +\Delta \mbox{ if } \delta > 0 \mbox{ else } -\Delta
  \theta_k^{(t+1)} := \theta_k^{(t)} + \mathrm{sign}_k\cdot\mathit{step}
  \theta_k^{(t+1)} := \min(\theta_k^{\max},\max(\theta_k^{\min},\theta_k^{(t+1)}))
|od|
|return| \bm{\theta}^{(t+1)},\;\bar\sigma
\end{program}
\hrule
\end{minipage}
\end{center}

\begin{table}[h]\centering
  \caption{Explore direction $\mathrm{sign}_k$ for target parameters.}
  \begin{tabular}{lll}
    \toprule Parameter & $\mathrm{sign}_k$ & Reasoning \\ \midrule
    $\rho$ (evaporation) & $+1$ & Faster evaporation forgets sooner $\to$ explore \\
    $\alpha$ (trail exponent) & $-1$ & Higher $\alpha$ trusts trail $\to$ exploit \\
    $\beta$ (visibility exp) & $-1$ & Higher $\beta$ is greedy $\to$ exploit \\
    $f$ (ACO fraction) & $-1$ & Shift budget away from colony $\to$ explore \\
    \bottomrule
  \end{tabular}
\end{table}

\subsection{Null Controller (Ablation)}

Returns $\bm{\theta}^{(t+1)}=\bm{\theta}^{(t)}$. Adaptive and frozen arms execute
identical code paths, differing only in controller logic, drawing from the same random
stream.

% =============================================================
\section{Evaluation Framework}
\label{sec:evaluation}
% =============================================================

\subsection{Recovery Metrics}

\begin{definition}[Disruption Transient]
  For a scoreable disruption at time $t$ (preceded by $C_b=2$ quiet steps):
  \begin{align}
    b &= \frac{1}{C_b}\sum_{\ell=t-C_b}^{t-1} \excess^{(\ell)} \\
    \text{damage} &= \max_{\ell\in[t,\,t+3]}\excess^{(\ell)} - b \\
    \text{recovery\_time} &= \min\{k\ge t \mid \excess^{(k)} \le b + 0.1|b|\}
  \end{align}
\end{definition}

\subsection{Statistical Framework}

\paragraph{Pairing Strategy.}
Comparisons are paired on \emph{trials}, not timesteps, because timesteps within a trial
share a graph and event sequence. Treating them as independent inflates the sample size.
Each trial is reduced to a mean excess $\bar\excess$.

\paragraph{Holm Correction.}
A signal-by-scenario grid produces $M$ simultaneous comparisons. Holm's step-down
procedure adjusts the raw Wilcoxon signed-rank p-values:
$p_{(k)}^{\mathrm{Holm}} = \min(1,\;(M-k+1)p_{(k)})$,
enforcing monotonic significance thresholds.

\paragraph{Effect Size (Cliff's $\delta$).}
Measures stochastic dominance $\hat\delta \in [-1,1]$. With $T=400$ trials, negligible
differences can yield tiny p-values. $\hat\delta$ provides scale: $<0.147$ is negligible.

% =============================================================
\section{Experimental Results}
\label{sec:results}
% =============================================================

\textbf{Setup:} Geometric graph ($n=200$). $B=300$ evals. $T=400$ trials.

\begin{table}[h]\centering
  \caption{Mean excess ($\excess$) over moving optimum.}
  \begin{tabular}{lllll}
    \toprule Arm & Static & Low & Moderate & High \\ \midrule
    A$^*$ & 0.0\% & 0.0\% & 0.0\% & 0.0\% \\
    ACO & 8.1\% & 8.0\% & 7.2\% & 8.2\% \\
    PSO & 38.4\% & 37.3\% & 37.0\% & 35.4\% \\
    Hybrid-Fixed & 4.1\% & 5.0\% & 5.1\% & 6.0\% \\
    Hybrid-Adaptive & 5.8\% & 5.0\% & 5.0\% & 5.1\% \\ \bottomrule
  \end{tabular}
\end{table}

\begin{table}[h]\centering
  \caption{Ablation: Hybrid-Adaptive vs Hybrid-Fixed.}
  \begin{tabular}{lllll}
    \toprule Scenario & $\Delta$ (Adapt $-$ Fixed) & $p^{\mathrm{Holm}}$ & $\hat\delta$ & Verdict \\ \midrule
    Static   & $+0.017$ & 0.003 & 0.07 & Fixed significantly better \\
    Low      & $-0.001$ & 0.71 & 0.02 & Not significant \\
    Moderate & $-0.001$ & 0.94 & 0.01 & Not significant \\
    High     & $-0.009$ & 0.23 & 0.04 & Not significant \\ \bottomrule
  \end{tabular}
\end{table}

% =============================================================
\section{Conclusion}
\label{sec:conclusion}
% =============================================================

We formulated four path-finding algorithms (A$^*$, MMAS, PSO, Hybrid) and an adaptive
controller. The key findings are:

\begin{enumerate}[label=(\arabic*)]
  \item \textbf{Parameter quality dominates.} Tuned fixed parameters reduce excess by
        $\approx 3.5$ percentage points over presets.
  \item \textbf{Online adaptation does not.} Adapting tuned parameters reduces excess
        by $\approx 0.004$ points, indistinguishable from noise at $n=400$.
  \item \textbf{Hybrid architecture works.} Hybrid arms significantly outperform plain
        ACO and PSO across all disruption scenarios.
\end{enumerate}

Future work will test the allocation thesis (recency-driven hybrid budget splitting) and
investigate reinforcement-learning controllers and neural metaheuristics.

\begin{thebibliography}{9}
\bibitem{stutzle2000} T.~Stützle and H.~H.~Hoos. \emph{MAX--MIN Ant System}. FGCS, 16(8):889--914, 2000.
\bibitem{mavrovouniotis2017} M.~Mavrovouniotis and S.~Yang. \emph{Adapting the Pheromone Evaporation Rate in Dynamic Environments}. EC, 2017.
\end{thebibliography}
\end{document}

